\chapter*{Abstract}
\addcontentsline{toc}{chapter}{Abstract}

This project presents the development and implementation of an automated Optical Mark Recognition (OMR) system designed to streamline the process of evaluating multiple-choice examinations and surveys. The OMR Checker system addresses the critical need for efficient, accurate, and cost-effective assessment tools in educational institutions and organizations.

The system leverages advanced computer vision techniques and machine learning algorithms to process scanned OMR sheets with high accuracy and speed. Key features include robust image preprocessing capabilities that handle various scanning conditions, automatic template alignment, and support for multiple OMR sheet formats. The system can process over 200 OMR sheets per minute while maintaining nearly 100\% accuracy on high-quality scans and approximately 90\% accuracy on mobile-captured images.

The implementation utilizes Python as the primary programming language, with OpenCV for image processing, NumPy for numerical computations, and pandas for data management. The system architecture follows a modular design pattern, enabling easy customization and extension for different OMR layouts and requirements.

Extensive testing was conducted using various sample datasets, including low-resolution scans, xeroxed sheets, and mobile-captured images. The results demonstrate the system's robustness and reliability across different input conditions. The system successfully handles challenges such as image rotation, varying lighting conditions, and different paper qualities.

The OMR Checker system provides significant advantages over traditional manual evaluation methods, including reduced processing time, elimination of human error, cost-effectiveness, and scalability for large-scale assessments. The open-source nature of the project promotes community collaboration and continuous improvement.

This project contributes to the field of educational technology by providing a practical, efficient solution for automated assessment processing. The system's flexibility and accuracy make it suitable for various applications, from small classroom quizzes to large-scale standardized testing programs.

\textbf{Keywords:} Optical Mark Recognition, Computer Vision, Educational Technology, Automated Assessment, Image Processing, Python, OpenCV